
We compare GraphWalker with a comprehensive set of baselines.
Following previous work, the Freebase dump is acquired from \url{https://developers.google.com/freebase?hl=en}. 
We deploy Freebase with Virtuoso. 
All the baselines are evaluated on Freebase database.


\textbf{IO-Prompt.}
IO-Prompt evaluates the intrinsic parametric knowledge of LLMs by directly feeding the natural language question and requesting a final answer. No external knowledge graph or intermediate reasoning steps are involved. This setting serves as a baseline for the model's zero-shot performance without grounding.


\textbf{Reasoning on Graphs (RoG).}
RoG \citep{luo2023reasoning} decouples KGQA into planning and retrieval-reasoning. It first fine-tunes an LLM to generate relation paths as reasoning plans, which are then used to retrieve relevant subgraphs from the KG. Finally, a reader module reasons over the retrieved subgraphs to answer the question. Unlike GraphWalker, RoG's planning is static and does not dynamically interact with the KG environment during the planning phase.

\textbf{Think-on-Graph (ToG).} 
ToG \citep{sun2023think} integrates structured knowledge from a knowledge graph into the LLM’s reasoning process, representing the methods using LLM generation to perform knowledge retrieval. It utilizes KG-based retrieval to enrich the model’s input, allowing it to leverage explicit entity-relation structures for improved factual accuracy. At each step, it uses the LLM to decide whether to continue retrieving or terminate the process.

\textbf{Think-on-Graph 2.0 (ToG-2.0).}
ToG-2 \citep{ma2025thinkongraph20deepfaithful} introduces a tightly-coupled hybrid RAG framework that iteratively integrates structured KGs and unstructured document contexts. It alternates between knowledge-guided text retrieval and context-enhanced graph search to gather in-depth heterogeneous clues. At each step, the LLM evaluates the retrieved knowledge to decide whether to continue exploring or generate the final answer.


 
\textbf{Generate-on-Graph (GoG).}
GoG \citep{xu2024generate} addresses KGQA under incomplete knowledge graphs via a training-free Thinking-Searching-Generating framework, where the LLM acts simultaneously as an agent to explore the KG and as a knowledge source to generate missing factual triples on demand. When retrieved subgraphs lack sufficient information, GoG invokes a Generate Action that synthesizes new triples grounded in the explored context and verifies them before reasoning, enabling the model to bridge structural gaps in the KG using its parametric knowledge.

\textbf{KG-Agent.}
KG-Agent \citep{jiang2025kg} adopts a tool-learning paradigm. It fine-tunes a smaller LLM (e.g., Llama-2-7B) to act as an autonomous agent that synthesizes executable tool calls (e.g., logical operations, KG queries). The model is trained on a synthesized dataset derived from functional programs, enabling it to decompose complex questions into sequences of tool executions.

\textbf{KBQA-o1.}
KBQA-o1 \citep{luo2025kbqa} proposes an agentic KBQA framework that combines a ReAct-based agent process with Monte Carlo Tree Search (MCTS) for heuristic KB environment exploration, where a policy model guides stepwise logical form generation and a reward model scores complete reasoning trajectories to balance search effectiveness and efficiency. To reduce reliance on human annotation, it further employs MCTS-driven exploration over unlabeled questions to auto-generate high-quality training data for incremental fine-tuning of both models.

\textbf{KG-R1.}
KG-R1~\citep{song2025efficient} is a state-of-the-art RL framework for KGQA that trains a lightweight LLM as a reasoning agent, optimizing its retrieval and reasoning policy end-to-end via turn-level rewards to navigate the KG efficiently. Unlike GraphWalker, KG-R1 does not employ a decoupled exploration pre-training stage. For a fair comparison, we equip KG-R1 with \textit{the same} four tools as specified in their original paper.